\documentclass{article}
\usepackage[utf8]{inputenc}
\usepackage[T1]{fontenc}
\usepackage{geometry}
\geometry{
    a4paper,
    left=2.5cm,
    right=2.5cm,
    top=2.5cm,
    bottom=2.5cm
}
\usepackage{graphicx}
\usepackage{float}
\usepackage{hyperref}
\usepackage[english]{babel}
\usepackage{amsmath}
\usepackage{amssymb}
\usepackage{array}
\usepackage{booktabs}
\usepackage{caption}
\usepackage{subcaption}
\usepackage{enumitem}
\usepackage{fancyhdr}
\usepackage{titlesec}
\usepackage{url}
\usepackage{textcomp}
\usepackage{xcolor}
\usepackage{listings}
\usepackage{multirow}
\usepackage{longtable}
\usepackage{tabularx}
\usepackage{times}

% Set up page style
\pagestyle{fancy}
\fancyhf{}
\rfoot{\thepage}
\renewcommand{\headrulewidth}{0pt}

% Title formatting
\titleformat{\chapter}[display]{\normalfont\bfseries}{\chaptertitle\thechapter}{0pt}{\Huge}
\titlespacing*{\chapter}{0pt}{50pt}{40pt}

% Document info
\title{Template Document}
\author{Academic Support Agent}
\date{\today}

\begin{document}
\maketitle

\pagestyle{empty}

\begin{minipage}{6cm} %Logo do IFPR
\resizebox{6cm}{1.6cm}{
\begin{tikzpicture}[square/.style={regular polygon,regular polygon sides=4}]
\tikzset{mynode/.style={square,rounded corners,draw=mygreen, top color=mygreen, bottom color=mygreen,very thick, text centered},}  
\node[circle, draw=myred , top color=myred, bottom color=myred] (q1) {$\quad$};
\node[mynode, below=0.08cm of q1] (q2) {$\;\;$};
\node[mynode, below=0.08cm of q2] (q3) {$\;\;$};
\node[mynode, below=0.08cm of q3] (q4) {$\;\;$};
\node[mynode, right=0.08cm of q1] (q5) {$\;\;$};
\node[mynode, below=0.08cm of q5] (q6) {$\;\;$};
\node[mynode, below=0.08cm of q6] (q7) {$\;\;$};
\node[mynode, below=0.08cm of q7] (q8) {$\;\;$};
\node[mynode, right=0.08cm of q5] (q9) {$\;\;$};
\node[mynode, right=0.08cm of q7] (q10) {$\;\;$};
\node[right=0.08cm of q8] (aux) {};
\node[text centered, right=0.06cm of q10] (q11) {\huge \sf \textbf{INSTITUTO FEDERAL}};
\node[mygreen, text centered, right=0.5cm of aux] (q12) {\huge\textsf{Paran }};
\end{tikzpicture}}
\end{minipage}
\begin{minipage}{7cm}
\begin{center}
\begin{footnotesize}
\sf \textbf{INSTITUTO FEDERAL DO PARAN }\\
\textbf{Campus Paranava }\\
Rua Jos  Felipe Tequinha, 1400\\
Jardim das Na es - Paranava  - PR
\end{footnotesize}
\end{center}
\end{minipage}
\begin{minipage}{4cm}
\centering
\includegraphics[scale=0.12]{meBrasil.png}\\
\begin{scriptsize}
\textsf{Minist rio da Educa o}
\end{scriptsize}
\end{minipage}
\doublespacing
\begin{center}
\begin{tabular}{|p{6.5cm}|p{6.5cm}|p{3.5cm}|}
\hline
\textbf{Curso:} \curso & \textbf{Turno:} \turno & \textbf{Trimestre:} \trimestre \\
\hline
\textbf{Disciplina:} \disciplina & \textbf{Professor(a):} \professor & \textbf{Data:} \\
\hline
\multicolumn{2}{|l|}{\textbf{Nome:}} & \textbf{R.A.:}\\
\hline
\textbf{Conceito:} $\square$ A $\;\;\square$ B $\;\;\square$ C $\;\;\square$ D & \multicolumn{2}{|l|}{\textbf{\instrumento}} \\
\hline
\end{tabular}
\end{center}
\singlespacing

%-------------------------------------------------------------------------
% EXPECTATIVAS DE APRENDIZAGEM 
%-------------------------------------------------------------------------

\begin{quote}
\begin{scriptsize}
\textbf{Expectativas de aprendizagem}: % (ou objetivos ou descri o dos crit rios de avalia o) 
 I -- Resolver corretamente e de forma completa os exerc cios; II -- Resolver os exerc cios no tempo definido; III -- Realizar a interpreta o correta dos exerc cios e IV -- Verificar o conhecimento adquirido em sala.\\
% Outro exemplo:
\textbf{Objetivos da avalia o:}:   VERIFICAR SE O ALUNO: -Identifica as principais fun es e rela es trigonom tricas;  identifica a aplicabilidade de teoria de seno, cosseno e tangente;  Identifica elementos do arco e opera com os mesmos;  Domina conceitos e c lculos no ciclo trigonom trico; Domina estrat gias de c lculos apropriados.\\
\textbf{Instru es}:  (Exemplo: Respostas   caneta nessa folha e com as resolu es em folhas anexas. Individual, sem consulta e sem uso de calculadora ou instrumentos eletr nicos. Dura o: duas aulas.)
\end{scriptsize}
\end{quote}

\hrule
\begin{center}
\begin{large}
\bf Avalia o de \disciplina \\
\end{large}
\end{center}
\hrule


%-------------------------------------------------------------------------
% QUEST ES ou ENUNCIADOS
%-------------------------------------------------------------------------

\begin{multicols}{2}  
% Caso n o queira que as quest es estejam dispostas em colunas basta comentar o comando acima e o comando \end{multicols}  no final do ambiente QUEST ES e ENUNCIADOS

\begin{enumerate} % Ambiente enumerado
\item Enunciado quest o 1

\item Enunciado quest o 2

\item Quest o com texto matem tico: O texto pode estar na linha do texto, como por exemplo: $\int x^3 ~ dx$ ou estar em linha pr pria: % clique sobre uma equa o

$$x=\frac{-b\pm \sqrt{b^2-4ac}}{2a}$$

\item Quest o com espa o em branco pr  estabelecido para resposta.

\vspace{3.5cm} % Escolha o tamanho do espa o em branco mudando o valor em cm

\item Quest o com figura: %   necess rio fazer upload da figura

\begin{center}
\includegraphics[scale=0.2]{meBrasil.png} \\ % O nome da figura deve ser id ntico ao nome do arquivo da figura.   poss vel alterar a escala
{\footnotesize Legenda da Figura} % Caso n o seja necess rio, apague esta linha
\end{center}

\item M ltipla escolha
% Alternativas
\begin{enumerate}
    \item Alternativa 1
    \item Alternativa 2
    \item Alternativa 3
    \item Alternativa 4
    \item Alternativa 5
\end{enumerate}

\item Verdadeiro ou falso?
% Afirma es
\begin{description}
    \item[(\hspace{0.5cm})] Afirma o 1
    \item[(\hspace{0.5cm})] Afirma o 2
    \item[(\hspace{0.5cm})] Afirma o 3
    \item[(\hspace{0.5cm})] Afirma o 4
    \item[(\hspace{0.5cm})] Afirma o 5
\end{description}

\item Quest o discursiva com linhas pontilhadas:
% Linhas: cada \dotfill abaixo representa uma linha pontilhada. Voc  pode acrescentar mais linhas ou tir -las.

\dotfill

\dotfill

\dotfill

\dotfill

\item Quest o discursiva com linhas cont nuas:
%caso deseje criar linhas cont nuas use \hrulefill

\hrulefill

\hrulefill

\hrulefill

\hrulefill

\end{enumerate}

\end{multicols}

%-------------------------------------------------------------------------
% MODELO DE CRIT RIOS que podem constar na avalia o
%-------------------------------------------------------------------------
\newpage % cria a quebra de p gina

Leia com aten o\footnote{Confirmo que li todas as informa es da atividade avaliativa com aten o e que tenho conhecimento dos crit rios de avalia o descritos acima.

\

Data: \underline{\hspace{0.7cm}} / \underline{\hspace{0.7cm}} / \underline{\hspace{0.7cm}} $\quad$ Nome: \underline{\hspace{11cm}}}:

\begin{center}
    \begin{tabular}{ | l | p{15cm} |}
\multicolumn{2}{l}{}\\    
\hline
Conceito A &  Quando 	a aprendizagem foi PLENA. Nesta avalia o significa que o aluno respondeu corretamente todas as quest es de forma completa e correta.\\
\hline
Conceito B & Quando 	a aprendizagem foi PARCIALMENTE PLENA. Nesta avalia o significa que o aluno respondeu correta e completamente 3 (tr s) quest es e 1 (uma) quest o parcialmente.\\
\hline
Conceito C     &  Quando 	a aprendizagem foi SUFICIENTE. Nesta avalia o significa que o aluno respondeu completa e corretamente 3 (tr s) quest es.\\
\hline
Conceito D     & Quando 	a aprendizagem foi INSUFICIENTE. Nesta avalia o significa que o aluno N O 	respondeu completa e corretamente no m nimo 3(tr s), n o atendendo de forma suficiente as expectativas de aprendizagem. \\
    \hline
    \end{tabular}
\end{center}
\end{document}